% §III --- System Model (~2 pages, ~1800 words). Drafted D6.
% All formulas track ``UavNetEnv.environments.py'' bit-exactly (the ground
% truth for the empirical results). Conference symbols are reused; new
% multi-UAV / dynamic terms are introduced where indicated.
\section{System Model}\label{sec:system}

We model a UAV-assisted multi-user MEC network in which a fleet of UAVs
collects visual streams from onboard cameras, encodes them into compact
DeepSC-VQA semantic features, and offloads the features through a base
station to an MEC server that performs cross-modal fusion with the
text queries streamed from the ground users. Time is slotted into
$T_s$ discrete slots of duration $S$, indexed by
$t\in\mathcal{T}=\{1,\dots,T_s\}$. Conference-version notation
\cite{xie2022taskoriented,liu2024jamminguav} is preserved; the
multi-UAV / dynamic-channel additions are introduced as needed.

\subsection{Network Model}\label{subsec:network}
The cell of radius $R_{\mathrm{cell}}$ contains a single base station
co-located with the MEC server at $d_{\mathrm{BS}}=(x_{\mathrm{BS}},
y_{\mathrm{BS}}, H_{\mathrm{BS}})$, $M$ UAVs indexed by
$\mathcal{U}=\{1,\dots,M\}$, and $N$ ground users indexed by
$\mathcal{N}=\{1,\dots,N\}$. We write $\mathcal{I}=\mathcal{U}\cup
\mathcal{N}$ for the set of all transmitting devices and use $|i|=M+N$.
UAV $m\in\mathcal{U}$ has position $d_m(t)=(x_m(t),y_m(t),H_U)$ at slot
$t$ with the altitude $H_U$ kept constant; UE $n\in\mathcal{N}$ has
position $d_n(t)=(x_n(t),y_n(t),H_n)$. The trajectory of every UAV is
constrained inside the cell to ensure reliable connectivity. UAV $m$
extracts visual semantic features from aerial imagery, while UE $n$
generates a textual query about the same scene; both feature streams are
forwarded through the BS to the MEC server, which fuses them and emits
the answer over a coordination control channel. Logical control signals
for task request and status feedback share a separate dedicated channel
and are not modeled in the resource allocation problem.

UEs form clusters with mobility within the cell. The number of UEs that
each UAV serves is decided online by a closest-UAV association rule:
\begin{equation}
\ell(n,t) \;=\; \arg\min_{m\in\mathcal{U}} \|d_m(t)-d_n(t)\|_2,
\qquad n\in\mathcal{N},
\label{eq:assoc}
\end{equation}
where $\ell(n,t)\in\mathcal{U}$ denotes the index of UAV that serves UE
$n$ at slot $t$. The association is recomputed at every slot.

\subsection{Channel Model}\label{subsec:channel}
The wireless channel between every device $i\in\mathcal{I}$ and the BS
follows a statistical air-to-ground propagation model
\cite{alhourani2014lap} that combines a line-of-sight (LoS) and a
non-line-of-sight (NLoS) component. The LoS probability is a logistic
function of the elevation angle $\theta_{t,i}$:
\begin{equation}
P^{\mathrm{LoS}}_{t,i}
\;=\; \frac{1}{1+\omega_{1}\exp\!\left\{-\omega_{2}(\theta_{t,i}-\omega_{1})\right\}},
\label{eq:plos}
\end{equation}
with environment-dependent constants $(\omega_{1},\omega_{2})$, and the
complementary NLoS probability $P^{\mathrm{NLoS}}_{t,i}=1-
P^{\mathrm{LoS}}_{t,i}$. Combining the path-loss exponents
$(\eta_{\mathrm{LoS}},\eta_{\mathrm{NLoS}})$ with a free-space term in
$f_c$, the large-scale channel gain in dB is
\begin{equation}
h^{\mathrm{ls}}_{t,i}
\;=\; \big(\eta_{\mathrm{LoS}}-\eta_{\mathrm{NLoS}}\big)P^{\mathrm{LoS}}_{t,i}
       + \eta_{\mathrm{NLoS}}
       + 20\log_{10}\!\frac{4\pi f_c\, d_{t,i}}{c} + \kappa_0,
\label{eq:hls}
\end{equation}
where $d_{t,i}=\|d_i(t)-d_{\mathrm{BS}}\|_2$, $c$ is the speed of
light, and $\kappa_0$ is a fixed offset ($25$~dB in our setup). Unlike
the conference single-UAV / static-channel formulation
\cite{liu2024jamminguav}, the journal extension adds a Rayleigh
small-scale fading component that is resampled every $K_{\mathrm{f}}$
slots. Let $|\mathbf{g}_{t,i}|^2$ denote the per-slot Rayleigh power
averaged over $N_r$ receive antennas; the effective per-device dB gain
becomes $h_{t,i}=h^{\mathrm{ls}}_{t,i}-10\log_{10}|\mathbf{g}_{t,i}|^2$.
With orthogonal in-band sharing, the uplink SNR of device $i$ at slot
$t$ on its assigned channel $b_i(t)\in\mathcal{C}=\{0,1\}$ is
\begin{equation}
\gamma_{t,i} \;=\; p_{t,i}\,h_{t,i}\big/(B_{t,i}\,N_{0}),
\label{eq:snr}
\end{equation}
where $p_{t,i}$ is the transmit power, $B_{t,i}=B_c/n_{t,b_i(t)}$ is the
effective bandwidth allocated to $i$ when $n_{t,b_i(t)}=
|\{j\in\mathcal{I}:b_j(t)=b_i(t)\}|$ devices share channel $b_i(t)$,
and $N_0$ is the noise power spectral density.

\subsection{Semantic Transmission Modeling}\label{subsec:sem}
Following the task-oriented multi-user semantic communication framework
of \cite{xie2022taskoriented}, the semantic similarity at slot $t$
between UAV $\ell(n,t)$ and UE $n$ is
\begin{equation}
\xi^{(\ell(n,t),n)}_{t}
\;=\; f\!\big(\gamma_{t,\ell(n,t)},\,\gamma_{t,n};\,K_{t,\ell(n,t)},K_{t,n}\big),
\label{eq:simlut}
\end{equation}
where $K_{t,i}$ denotes the cardinality of the semantic-symbol set used
by device $i$. The mapping $f(\cdot)$ is implemented through a look-up
table (LUT) trained offline on the DeepSC-VQA model
\cite{xie2022taskoriented} and stored in the env package; the LUT
discretizes both SNR axes onto $[-10\!:\!5\!:\!20]$~dB bins. We write
$R_{t,i}$ for the resulting semantic transmission rate of node $i$:
\begin{equation}
R_{t,i} \;=\; \xi^{(\ell(n,t),n)}_t
       \,\frac{B_{t,i}}{K_{t,i}}\,\mathcal{H}_{i},
\label{eq:semrate}
\end{equation}
where $\mathcal{H}_{i}\in\{\mathcal{H}_{\mathrm{img}},\mathcal{H}_{\mathrm{txt}}\}$
is the semantic entropy of the image or text feature carried by device
$i$. Equation~\eqref{eq:semrate} provides a dimensionally consistent
metric (semantic units per second) for evaluating end-to-end information
delivery in the semantic domain.

\subsection{Task Arrival and UE Mobility}\label{subsec:dynamics}
The conference formulation assumed a fixed per-UE task workload and
stationary UEs, but realistic deployments inject \emph{both} traffic
randomness and topology drift. We capture both with two independent
processes that are switched on by the policy in dynamic-mode evaluation
(Section~\ref{sec:exp}).

\paragraph{Poisson task arrival.} At every slot, each UE $n$ samples a
new image-task count and a new text-task count from
$Q^{\mathrm{img}}_{n}(t),Q^{\mathrm{txt}}_{n}(t)\sim
\mathrm{Poisson}(\bar{\lambda})$, independently across slots and across
UEs. The semantic-symbol payload of UE $n$ at slot $t$ is therefore
$Q^{\mathrm{img}}_{n}(t)\,\mathcal{H}_{\mathrm{img}}+
Q^{\mathrm{txt}}_{n}(t)\,\mathcal{H}_{\mathrm{txt}}$, scaled by the
device-side task size constants. Conference behavior is recovered by
freezing the per-UE counts at the values sampled at the episode
beginning.

\paragraph{Slow random walk.} Each UE performs an independent bounded
2-D random walk in the horizontal plane: at slot $t$ the UE displacement
is $\Delta d_{n}(t)=(v_{n}(t)\cos\phi_{n}(t),v_{n}(t)\sin\phi_{n}(t))$
with speed $v_{n}(t)\sim\mathcal{U}[0,v_{\max}]$ and direction
$\phi_{n}(t)\sim\mathcal{U}[0,2\pi)$. Positions are clipped back into
the cell, and the dynamic-mode flag $v_{\max}>0$ controls activation.

\subsection{Delay and Energy Consumption Model}\label{subsec:cost}
The total task delay at slot $t$ is the sum of the semantic
transmission delay $T^{\mathrm{tr}}(t)$ to the MEC server and the
computation delay $T^{\mathrm{cmp}}_{\mathrm{MEC}}(t)$ for VQA
inference. With $T_{\mathrm{avail}}(t)=(T_s-t+1)S$ as the residual time
budget at slot $t$, the per-UAV transmission delay aggregates contributions
from the UEs it currently serves:
\begin{equation}
T^{\mathrm{tr}}_{\!U,m}(t)
\;=\; \frac{\sum_{n:\ell(n,t)=m} Q^{\mathrm{img}}_{n}(t)\mathcal{H}_{\mathrm{img}}}
            {R_{t,m}},
\quad m\in\mathcal{U},
\label{eq:tu}
\end{equation}
and the per-UE transmission delay is
$T^{\mathrm{tr}}_{\!n}(t)=
\big(Q^{\mathrm{txt}}_{n}(t)\mathcal{H}_{\mathrm{txt}}\big)/R_{t,n}$.
Because UAV and UE uplinks proceed in parallel over orthogonal
channels, the effective per-slot transmission delay is
\begin{equation}
T^{\mathrm{tr}}(t) \;=\;
\min\!\Big\{\,
   \max\!\big(\max_{m}T^{\mathrm{tr}}_{\!U,m}(t),\,
              \max_{n}T^{\mathrm{tr}}_{\!n}(t)\big),\,
   T_{\mathrm{avail}}(t)\Big\}.
\label{eq:ttotal}
\end{equation}
The MEC server processes the aggregated semantic data at constant
throughput $f_{\mathrm{MEC}}$, giving
$T^{\mathrm{cmp}}_{\mathrm{MEC}}(t)=\mathcal{A}_{\mathrm{MEC}}/
f_{\mathrm{MEC}}$ where $\mathcal{A}_{\mathrm{MEC}}=
\sum_{n}\!Q^{\mathrm{img}}_{n}(t)\mathcal{H}_{\mathrm{img}}+
\sum_{n}\!Q^{\mathrm{txt}}_{n}(t)\mathcal{H}_{\mathrm{txt}}$. Total
slot delay is $T^{\mathrm{tot}}(t)=T^{\mathrm{tr}}(t)+
T^{\mathrm{cmp}}_{\mathrm{MEC}}(t)$, subject to the scheduling
constraint $T^{\mathrm{tot}}(t)\le T_{\mathrm{avail}}(t)$.

The system energy consumption at slot $t$ aggregates UAV mobility,
hover-stage transmission, and per-UE transmission. The mobility energy of
UAV $m$ is $E^{\mathrm{mob}}_{U,m}(t)=P_{\mathrm{fly}}\,t\,S+
P_{\mathrm{hov}}\,T^{\mathrm{tr}}_{\!U,m}(t)$ where
$P_{\mathrm{fly}}$ and $P_{\mathrm{hov}}$ are the propulsion and hover
power. The transmission energy is
$E^{\mathrm{tr}}_{i}(t)=P_{i}(t)T^{\mathrm{tr}}_{i}(t)$ for every
device. Letting $E^{\max}$ denote the maximum admissible per-slot
energy under worst-case full-load operation, the joint system
consumption index is
\begin{equation}
O(t) \;=\; \alpha\,\widetilde{T}(t)
       + (1-\alpha)\,\widetilde{E}(t),
\quad
\widetilde{T}(t)=\frac{T^{\mathrm{tot}}(t)}{T_{\mathrm{avail}}(t)},
\;\;
\widetilde{E}(t)=\frac{E^{\mathrm{sys}}(t)}{E^{\max}},
\label{eq:cost}
\end{equation}
both normalized to $[0,1]$. The trade-off coefficient $\alpha\in[0,1]$
encodes mission preferences (e.g., a delay-critical mission selects a
larger $\alpha$); we treat $\alpha$ as an external knob and study the
sensitivity of the learned policy to its choice in
Section~\ref{sec:exp:sens}. The full optimization problem and Dec-POMDP
casting follow in Section~\ref{sec:method}.
